\documentclass{report}
\usepackage{amsmath,amssymb}
\setlength{\parindent}{0mm}
\setlength{\parskip}{1em}
\begin{document}
\begin{center}
\rule{6in}{1pt} \
{\large Charles Morgenstern \\
{\bf Efficient Simulation of High-Frequency Wave Propagation Using Domain Decomposition and High-Order FEM}}

Applied Mathematics \& Statistics \\ 1500 Illinois St \\ Golden \\ Colorado 80401
\\
{\tt cmorgens@mymail.mines.edu}\\
Mahadevan Ganesh\end{center}

Large scale scientific computing models, requiring iterative algebraic solvers,
are needed to simulate high-frequency wave propagation.
This is because large degrees of freedom are needed to avoid the
celebrated Helmholtz computer model pollution effects.
Using low-order finite difference or finite element methods (FDM/FEM),
such issues have been well investigated for low and medium frequency
models (typically at most $50$ wavelengths per diameter of the wave
propagation domain).
Standard FDM/FEM based discretizations of the time-harmonic Helmholtz
wave propagation model lead to sign-indefinite systems with eigenvalues
in the left half of the complex plane.
Hence standard iterative methods (such as GMRES/BiCGstab) perform poorly,
and additional techniques such as multigrid (MG) or decomposition of the
domain are required for efficient and practical simulation of
high-frequency FDM/FEM Helmholtz models.
In this work, we investigate the use of multiple additive Schwarz type
domain decomposition (DD) approximations to efficiently simulate
high-frequency wave propagation with high-order FEM.
We compare our DD based results with those obtained using a standard
geometric MG approach for over $100$ wavelength models.


\end{document}
